% =====================================================================
%  תרגיל בית 3 --- אחרי הרצאות 3--5 (הפרעות, הילות וחומר אפל, CMB)
%  קומפילציה: lualatex hw3.tex
% =====================================================================
\documentclass[11pt,a4paper]{article}
% =====================================================================
%  Shared preamble — lecture notes "Introduction to Cosmology" (9 h)
%  Compile:  lualatex cosmo_lectures.tex   (twice, for the TOC)
% =====================================================================
\usepackage[a4paper,top=2.2cm,bottom=2.2cm,left=2.2cm,right=2.2cm]{geometry}
\usepackage{fontspec}
\usepackage{amsmath,amssymb}
\usepackage[bidi=basic]{babel}
\babelprovide[import,main]{hebrew}
\babelprovide[import]{english}

% ---- fonts: Culmus (Overleaf/Linux) -> Windows David -> FreeSerif ----
\IfFontExistsTF{David CLM}{%
  \babelfont{rm}{David CLM}\babelfont{sf}{Simple CLM}\babelfont{tt}{Miriam Mono CLM}%
}{\IfFontExistsTF{David}{%
  \babelfont{rm}[ItalicFeatures={FakeSlant=0.15}]{David}\babelfont{sf}{Arial}\babelfont{tt}{Courier New}%
}{\IfFileExists{/mnt/c/Windows/Fonts/david.ttf}{%
  \babelfont{rm}[Path=/mnt/c/Windows/Fonts/,Extension=.ttf,BoldFont=davidbd,ItalicFont=david,
    ItalicFeatures={FakeSlant=0.15}]{david}%
  \babelfont{sf}[Path=/mnt/c/Windows/Fonts/,Extension=.ttf,BoldFont=arialbd]{arial}%
  \babelfont{tt}[Path=/mnt/c/Windows/Fonts/,Extension=.ttf]{cour}%
}{%
  \babelfont{rm}{FreeSerif}\babelfont{sf}{FreeSans}\babelfont{tt}{FreeMono}%
}}}
\babelfont[english]{rm}{Latin Modern Roman}
\babelfont[english]{tt}{Latin Modern Mono}

\usepackage{xcolor}
\usepackage[framemethod=default]{mdframed}
\usepackage{enumitem}
\usepackage{booktabs}
\usepackage{graphicx}
\usepackage{fancyhdr}
\usepackage{needspace}
\usepackage{listings}
\usepackage[hidelinks]{hyperref}
\lstset{basicstyle=\ttfamily\footnotesize,keywordstyle=\color{qblue}\bfseries,
  commentstyle=\color{ggreen}\itshape,frame=single,rulecolor=\color{black!30},
  backgroundcolor=\color{black!2},breaklines=true,showstringspaces=false,
  language=Python,columns=fullflexible}

\renewcommand{\labelitemi}{\ensuremath{\bullet}}
\renewcommand{\labelitemii}{\ensuremath{\circ}}
\setlist[itemize]{leftmargin=1.4em,itemsep=1pt,topsep=2pt}
\setlist[enumerate]{leftmargin=1.6em,itemsep=1pt,topsep=2pt}

% ---- colours (as in the original course) ----
\definecolor{qblue}{RGB}{20,60,120}
\definecolor{ggreen}{RGB}{20,100,70}
\definecolor{kpurple}{RGB}{95,45,110}
\definecolor{worange}{RGB}{170,90,10}

\newmdenv[linecolor=kpurple,linewidth=1pt,backgroundcolor=kpurple!4,
  innertopmargin=6pt,innerbottommargin=6pt,innerleftmargin=8pt,innerrightmargin=8pt,
  skipabove=7pt,skipbelow=7pt]{keybox}
\newmdenv[linecolor=worange,linewidth=1pt,backgroundcolor=worange!5,
  innertopmargin=6pt,innerbottommargin=6pt,innerleftmargin=8pt,innerrightmargin=8pt,
  skipabove=7pt,skipbelow=7pt]{warnbox}
\newmdenv[linecolor=ggreen,linewidth=1pt,backgroundcolor=ggreen!5,
  innertopmargin=6pt,innerbottommargin=6pt,innerleftmargin=8pt,innerrightmargin=8pt,
  skipabove=7pt,skipbelow=7pt]{goalbox}
\newmdenv[linecolor=qblue,linewidth=1.1pt,backgroundcolor=qblue!5,
  innertopmargin=6pt,innerbottommargin=6pt,innerleftmargin=8pt,innerrightmargin=8pt,
  skipabove=8pt,skipbelow=6pt]{exbox}
\newmdenv[linecolor=black!35,linewidth=0.7pt,backgroundcolor=black!3,
  innertopmargin=6pt,innerbottommargin=6pt,innerleftmargin=8pt,innerrightmargin=8pt,
  skipabove=4pt,skipbelow=8pt]{solbox}

% ---- macros ----
\newcommand{\en}[1]{\foreignlanguage{english}{#1}}
\newcommand{\Msun}{M_{\odot}}
\newcommand{\Mpc}{\mathrm{Mpc}}
\newcommand{\kpc}{\mathrm{kpc}}
\newcommand{\Gyr}{\mathrm{Gyr}}
\newcommand{\kms}{\mathrm{km\,s^{-1}}}
\newcommand{\kmsMpc}{\mathrm{km\,s^{-1}\,Mpc^{-1}}}
\newcommand{\LCDM}{\Lambda\mathrm{CDM}}

% lecture header: \lecture{number}{title}{time budget + one-line content}
\newcounter{lec}
\counterwithin*{subsection}{lec}
\renewcommand{\thesubsection}{\thelec.\arabic{subsection}}
\makeatletter\renewcommand{\fps@figure}{htbp}\makeatother
\newcommand{\lecture}[3]{%
  \clearpage\refstepcounter{lec}%
  \noindent\colorbox{qblue}{\parbox{\dimexpr\textwidth-2\fboxsep}{%
   \color{white}\bfseries\Large הרצאה #1: #2\\[3pt]\normalsize\mdseries #3}}%
  \addcontentsline{toc}{section}{הרצאה #1: #2}\par\medskip}

% exercises: numbered per lecture (1.1, 1.2, ...)
\newcounter{ex}[lec]
\newenvironment{exercise}[1]{%
  \refstepcounter{ex}\begin{exbox}\textbf{\textcolor{qblue}{תרגיל \thelec.\theex\ --- #1.}}\ }%
  {\end{exbox}}
\newcommand{\soln}[1]{\par\smallskip\noindent\textbf{פתרון #1.}\ }

\pagestyle{fancy}\fancyhf{}
\fancyhead[R]{\small מבוא לקוסמולוגיה --- רשימות הרצאה}
\fancyhead[L]{\small\hbox{\textdir TLT\relax\thepage}}
\renewcommand{\headrulewidth}{0.4pt}
\setlength{\parindent}{0pt}\setlength{\parskip}{4pt}

% ---- shared definitions for the homework sheets (input after ../lecture_notes/preamble) ----
\newcounter{hwq}
\newenvironment{question}[2]{%
  \refstepcounter{hwq}\needspace{5\baselineskip}%
  \begin{exbox}\textbf{\textcolor{qblue}{שאלה \thehwq\ --- #1}}\hfill\textbf{(#2 נק')}\par\smallskip}%
  {\end{exbox}}
\newenvironment{parts}{\begin{enumerate}[label=(\alph*),itemsep=3pt]}{\end{enumerate}}
\newcommand{\hwtitle}[3]{%
  \noindent\colorbox{qblue}{\parbox{\dimexpr\textwidth-2\fboxsep}{%
   \color{white}\bfseries\Large תרגיל בית #1: #2\\[3pt]\normalsize\mdseries #3}}\par\medskip}
\newcommand{\hwrules}[1]{%
\begin{goalbox}\small\textbf{הנחיות.}
\begin{itemize}
\item ארבע שאלות, 25 נקודות כל אחת. השאלות בלתי תלויות זו בזו.
\item אלא אם נאמר אחרת: $\LCDM$ שטוח עם $h=0.674$, $\Omega_m=0.315$, $\Omega_\Lambda=1-\Omega_m$ (קרינה מוזנחת), ו-$c/H_0=4448\,\Mpc$, $1/H_0=14.51\,\Gyr$.
\item רוב הסעיפים אנליטיים: גזירה, ואז הצבת מספרים במחשבון. רק סעיפים המסומנים \textbf{[מחשב]} דורשים קוד; בהם צרפו את הקוד (\en{Python} מומלץ), סמנו צירים ויחידות, ושמרו את ה-\en{seed}.
\item כל טענה, נתון או ערך שלא נלמדו בהרצאה --- עם ציטוט (מחבר, שנה, כתב עת).
#1
\end{itemize}\end{goalbox}}
\pagestyle{fancy}\fancyhf{}
\fancyhead[L]{\small\hbox{\textdir TLT\relax\thepage}}

\fancyhead[R]{\small מבוא לקוסמולוגיה --- תרגיל בית 3}
\begin{document}
\hwtitle{3}{הפרעות, הילות של חומר אפל ו-\en{CMB}}{אחרי הרצאות 3--5 \quad$\cdot$\quad פתרונות למשוואת הצמיחה, פרופיל \en{NFW}, עקומת סיבוב, קרינת הרקע}
\hwrules{\item $G=4.30\times10^{-6}\,\kpc\,(\kms)^2/\Msun$; $H_0=0.0674\,\kms\,\kpc^{-1}$.
\item הנתונים לשאלה 3 בטבלה בסוף התרגיל.}

% ---------------------------------------------------------------------
\begin{question}{פתרונות של משוואת הצמיחה הלינארית}{25}
משוואת ההפרעות לחומר קר (מהרצאה 3): $\ddot\delta+2H\dot\delta=4\pi G\bar\rho_m\delta$.
\begin{parts}
\item \textbf{(5)} \emph{עידן הקרינה.} כש-$\rho_r\gg\rho_m$, $H=1/2t$ והאיבר $4\pi G\bar\rho_m\delta$ זניח. פתרו והראו ש-$\delta=C_1+C_2\ln t$. למה הפרעות בחומר האפל כמעט לא גדלות לפני השוויון (אפקט \en{Mészáros})?
\item \textbf{(7)} עם $y=a/a_{\rm eq}$ ושני הרכיבים, המשוואה הופכת ל-
\[\frac{d^2\delta}{dy^2}+\frac{2+3y}{2y(1+y)}\frac{d\delta}{dy}-\frac{3}{2y(1+y)}\delta=0.\]
בדקו ש-$\delta\propto1+\tfrac32y$ פתרון. מהם הגבולות $y\ll1$ ו-$y\gg1$? בכמה גדלה הפרעה מ-$z=10^5$ ועד השוויון ($z_{\rm eq}=3400$), ובכמה גדל באותו זמן $a$?
\item \textbf{(8)} \emph{פתרון מדויק ליקום של חומר${}+\Lambda+{}$עקמומיות.} (i) השתמשו במשוואות פרידמן ($\dot H=-4\pi G\bar\rho_m+kc^2/a^2$) והראו ש-$\delta\propto H(t)$ פותר את משוואת הצמיחה. האם זה מוד גדל או דועך? (ii) הציבו $\delta=H(t)\,u(t)$ והראו ש-$\dot u\propto1/(aH)^2$, ולכן המוד הגדל הוא
\[D(a)\propto H(a)\int_0^a\frac{da'}{[a'H(a')]^3}.\]
(iii) בדקו שב-\en{EdS} מתקבל $D\propto a$. (iv) הראו שכאשר $\Lambda$ שולט ($H\to$ קבוע) $D\to$ קבוע. למה ביקום פתוח בלי $\Lambda$ הצמיחה נעצרת מוקדם יותר מאשר ב-$\LCDM$ עם אותו $\Omega_m$? (רמז: מתי $\Omega_k a^{-2}$ ו-$\Omega_\Lambda$ משתווים לחומר?)
\item \textbf{(5)} ב-$\LCDM$ האינטגרל מ-(ג) נותן $D(a{=}1)=0.79$ בנרמול $D\to a$ בעידן החומר. בכמה גדלה הפרעה מ-$z=1090$ עד היום? ב-\en{CMB} רואים $\delta T/T\sim10^{-5}$, כלומר הפרעות בריוניות $\delta_b\sim10^{-5}$ ברקומבינציה. האם בלי חומר אפל היו היום גלקסיות? מה חומר אפל קר משנה (רמז: הרצאה 3 ותרגיל מסת ג'ינס)?
\end{parts}
\end{question}

% ---------------------------------------------------------------------
\begin{question}{הילת \en{NFW} של שביל החלב}{25}
$\rho(r)=\dfrac{\rho_s}{(r/r_s)(1+r/r_s)^2}$, $M(<r)=4\pi\rho_sr_s^3\,\mu(r/r_s)$ עם $\mu(x)=\ln(1+x)-\frac{x}{1+x}$. $r_{200}$ מוגדר כרדיוס שבתוכו הצפיפות הממוצעת היא $200\rho_c$, ו-$c=r_{200}/r_s$.
\begin{parts}
\item \textbf{(6)} הראו ש-$r_{200}=\big(GM_{200}/100H_0^2\big)^{1/3}$, ש-$v_{200}\equiv\sqrt{GM_{200}/r_{200}}=10H_0r_{200}$, ושהצפיפות האופיינית היא $\rho_s/\rho_c=\frac{200}{3}\frac{c^3}{\mu(c)}$.
\item \textbf{(8)} להילה של שביל החלב קחו $M_{200}=1.2\times10^{12}\Msun$, $c=10$. חשבו $r_{200}$, $v_{200}$, $r_s$, $\rho_s$ (ב-$\Msun\,\kpc^{-3}$), ואת $v_{\max}$ ו-$r_{\max}$ (השתמשו ב-$x_{\max}=2.163$ מהתרגיל בהרצאה). מה מסת החומר האפל בתוך מסלול השמש ($8\,\kpc$) ומה תרומתה ל-$v_c$ שם? השוו ל-$v_c\approx230\,\kms$ הנמדד: מי עוד תורם?
\item \textbf{(5)} הראו ש-$M(<r)$ של \en{NFW} מתבדרת לוגריתמית. פי כמה גדולה המסה בתוך $2r_{200}$ ובתוך $10r_{200}$ מ-$M_{200}$? למה בכל זאת מגדירים מסה "ויריאלית" (רמז: הרצאה 4, $18\pi^2$)?
\item \textbf{(6)} \emph{קאספ מול ליבה.} מצאו את השיפוע הלוגריתמי $d\ln\rho/d\ln r$ של \en{NFW} כפונקציה של $x$, ואת גבולותיו. השוו לפרופיל \en{Burkert} $\rho\propto[(1+x)(1+x^2)]^{-1}$ ולספירה איזותרמית $\rho=\sigma^2/(2\pi Gr^2)$ (מה $v_c$ שלה?). איזה מבחן תצפיתי מבחין בין קאספ לליבה, ולמה הוא קשה בגלקסיות ננסיות?
\end{parts}
\end{question}

% ---------------------------------------------------------------------
\begin{question}{חומר אפל מעקומת סיבוב \textbf{[מחשב]}}{25}
בטבלה בסוף: עקומת סיבוב של גלקסיה מודל (\emph{נתונים סינתטיים}): $v_{\rm obs}\pm\sigma_v$, תרומת הדיסקה הכוכבית $v_{\rm disk}$ עבור יחס מסה-אור $\Upsilon_d=1$, ותרומת הגז $v_{\rm gas}$. המודל: $v^2=\Upsilon_dv_{\rm disk}^2+v_{\rm gas}^2+v_{\rm DM}^2$.
\begin{parts}
\item \textbf{(6)} שרטטו את הנתונים עם התרומות הבריוניות. התאימו בריונים בלבד ($v_{\rm DM}=0$, $\Upsilon_d$ חופשי): מה $\Upsilon_d$ המיטבי ו-$\chi^2/\nu$? באיזה רדיוס הכישלון בולט?
\item \textbf{(8)} הוסיפו הילת \en{NFW} (שאלה 2) וקבעו $\Upsilon_d=0.5$ (ערך טיפוסי ב-$3.6\,\mu$m). מצאו את $(\log_{10}M_{200},c)$ המיטביים על גריד או ב-\en{MCMC} מתרגיל 2, ושרטטו את קווי $\Delta\chi^2=2.30$ ו-$6.17$. מה $\chi^2/\nu$?
\item \textbf{(5)} מה יחס המסה האפלה לבריונית בתוך $30\,\kpc$? (השתמשו במודל המיטבי, ו-$M(<R)=v^2R/G$ לכל רכיב.) השוו ליחס הקוסמי $\Omega_{\rm dm}/\Omega_b\approx5.4$.
\item \textbf{(6)} \emph{ניוון דיסקה--הילה.} חזרו על (ב) עם $\Upsilon_d=0.3$ ו-$0.7$. איך משתנים $M_{200}$, $c$ ו-$\chi^2$? למה עקומת הסיבוב לבדה מתקשה לקבוע את $\Upsilon_d$, ואיזה מידע נוסף שובר את הניוון? (שאלת בונוס, 3 נק': התאימו גלקסיה אמיתית מ-\en{SPARC}, \en{Lelli, McGaugh \& Schombert 2016, AJ 152, 157}.)
\end{parts}
\end{question}

% ---------------------------------------------------------------------
\begin{question}{קרינת הרקע: פוטונים, דיפול ובעיית האופק}{25}
\begin{parts}
\item \textbf{(6)} מספקטרום גוף שחור עם $T_0=2.7255\,\mathrm K$ חשבו את צפיפות הפוטונים $n_\gamma=\frac{2\zeta(3)}{\pi^2}\big(\frac{k_BT_0}{\hbar c}\big)^3$ ב-$\mathrm{cm^{-3}}$ ואת היחס בריונים לפוטונים $\eta=n_b/n_\gamma$ (עם $n_b$ משאלה 1(ד) של תרגיל 1). הראו ש-$\eta$ קבוע בזמן מאז $e^\pm$ השמדה, ושהספקטרום נשאר של גוף שחור עם $T\propto1/a$.
\item \textbf{(5)} \emph{הדיפול.} הראו שצופה הנע במהירות $v\ll c$ יחסית לרקע רואה $T(\theta)\simeq T_0(1+\frac vc\cos\theta)$. הדיפול הנמדד הוא $3.36\,\mathrm{mK}$: מה המהירות? למה מורידים אותו לפני שמנתחים את האניזוטרופיות?
\item \textbf{(6)} \emph{\en{Sachs--Wolfe}.} בסקאלות גדולות $\delta T/T=\Phi/3c^2$. אם $\delta T\approx30\,\mu\mathrm K$, מה $\Phi/c^2$? השוו לפוטנציאל של צביר גלקסיות היום, $\Phi/c^2\sim\sigma^2/c^2$ עם $\sigma=1000\,\kms$. מה המשמעות: האם הפוטנציאלים גדלים בעידן החומר? (רמז: $\Phi\propto\delta/a$ בתרגיל הצמיחה.)
\item \textbf{(8)} \emph{בעיית האופק.} האופק החלקיקי הנע ברקומבינציה הוא $\eta_\star=\int_0^{t_\star}c\,dt/a$. (i) בקירוב של חומר בלבד הראו ש-$\eta_\star=\frac{2c}{H_0\sqrt{\Omega_m}}(1+z_\star)^{-1/2}$, ואז עם חומר${}+{}$קרינה (כמו אופק הקול בהרצאה 5, אבל עם $c$ במקום $c_s$). (ii) עם $D_M(z_\star)=13.9\,\mathrm{Gpc}$ ו-$z_\star=1090$, מה הגודל הזוויתי של אזור מקושר סיבתית בשמיים? (iii) לכמה כתמים בלתי תלויים מתחלקים השמיים? ובכל זאת $T$ אחידה עד $10^{-5}$ --- הסבירו את הבעיה במשפט, ואת הפתרון שיוצג בהרצאה 6.
\end{parts}
\end{question}

\clearpage
\begin{center}
\textbf{טבלה: עקומת סיבוב של גלקסיה מודל (לשאלה 3). מהירויות ב-$\kms$.}\\[6pt]
\begin{otherlanguage}{english}
\begin{tabular}{rrrrr}
\toprule
$R$ [kpc] & $v_{\rm obs}$ & $\sigma_v$ & $v_{\rm disk}\,(\Upsilon_d{=}1)$ & $v_{\rm gas}$\\
\midrule
1 & 79.9 & 5.8 & 107.7 & 4.3 \\
2 & 138.0 & 7.1 & 165.0 & 7.5 \\
3 & 161.9 & 8.0 & 199.6 & 10.0 \\
4 & 174.8 & 8.5 & 220.0 & 12.0 \\
5 & 174.1 & 8.9 & 230.8 & 13.6 \\
6 & 209.7 & 9.1 & 235.2 & 15.0 \\
8 & 201.3 & 9.3 & 231.9 & 17.1 \\
10 & 210.7 & 9.2 & 220.6 & 18.8 \\
12 & 212.6 & 9.2 & 206.4 & 20.0 \\
14 & 206.2 & 9.1 & 192.1 & 21.0 \\
16 & 199.5 & 9.0 & 178.8 & 21.8 \\
18 & 205.3 & 9.0 & 167.0 & 22.5 \\
20 & 202.7 & 8.9 & 156.7 & 23.1 \\
22 & 202.4 & 8.9 & 147.9 & 23.6 \\
24 & 181.6 & 8.9 & 140.3 & 24.0 \\
26 & 192.3 & 8.9 & 133.8 & 24.4 \\
28 & 189.4 & 8.9 & 128.2 & 24.7 \\
30 & 221.2 & 8.9 & 123.2 & 25.0 \\
\bottomrule
\end{tabular}
\end{otherlanguage}
\end{center}

\end{document}
